\section{DiT、Text-to-Image 和 Text-to-Video}

\subsection{为什么 Transformer 又回来了}

当 diffusion 进入 latent space 后，主干网络就越来越像一个序列建模问题。于是 Transformer 再次变得非常自然：token 序列、时间步条件、文本条件，都可以比较统一地接进去。DiT 的关键问题不是“能不能用 Transformer”，而是“怎么把条件信号优雅地注入进去”。

\begin{figure}[H]
\centering
\includegraphics[width=0.95\textwidth]{slides-images/slide-097.jpg}
\caption{Diffusion Transformer：把 diffusion 变成标准 Transformer block 的问题，核心是条件注入方式\protect\footnotemark}
\end{figure}
\footnotetext{Slides 第98页。}

\begin{figure}[H]
\centering
\includegraphics[width=0.95\textwidth]{slides-images/slide-099.jpg}
\caption{最常见的时间步注入方式之一：对归一化后的特征做 scale/shift 调制\protect\footnotemark}
\end{figure}
\footnotetext{Slides 第100页。}

\begin{knowledgebox}{DiT 里最重要的两类 conditioning}
\begin{itemize}
    \item \textbf{timestep conditioning}：通常用 scale-shift / AdaLN 方式注入
    \item \textbf{text or multimodal conditioning}：通常用 cross-attention 或 joint attention
\end{itemize}
\end{knowledgebox}

\subsection{Text-to-Image}

文本到图像的标准管线非常清楚：prompt 先进入预训练 text encoder，然后 text embedding 和 noisy latent 一起进入 diffusion transformer，最后经过 decoder 还原成图像。这里最关键的不是“把文字转成图片”这么一句话，而是“让条件信号真的参与了生成轨迹的每一步修正”。

\begin{figure}[H]
\centering
\includegraphics[width=0.95\textwidth]{slides-images/slide-100.jpg}
\caption{Text-to-image 的标准结构：text encoder、diffusion transformer 和 decoder 组成一条完整的条件生成链路\protect\footnotemark}
\end{figure}
\footnotetext{Slides 第101页。}

\begin{figure}[H]
\centering
\includegraphics[width=0.95\textwidth]{slides-images/slide-101.jpg}
\caption{一个工业级例子：FLUX.1 把文本编码、latent diffusion 和大规模 Transformer 组合到一起\protect\footnotemark}
\end{figure}
\footnotetext{Slides 第102页。}

\begin{importantbox}{为什么 text-to-image 看起来像“多模型拼装”}
因为它本来就是。文本编码器通常是预训练并冻结的；latent encoder / decoder 负责压缩和还原；真正学生成过程的是 diffusion 主干。工程上，每个模块都在做自己最擅长的事。
\end{importantbox}

\subsection{Text-to-Video}

视频生成比图像生成贵得多，原因不是“多了一维”这么简单，而是时序一致性、运动一致性、跨帧语义一致性都会把 token 数和计算成本推高。到了视频场景，latent 往往是 $t \times h \times w \times c$，模型要同时处理空间和时间。

\begin{figure}[H]
\centering
\includegraphics[width=0.95\textwidth]{slides-images/slide-102.jpg}
\caption{Text-to-video 的结构：同样是 text encoder + diffusion transformer + decoder，但 latent 里多了时间维度\protect\footnotemark}
\end{figure}
\footnotetext{Slides 第103页。}

\begin{figure}[H]
\centering
\includegraphics[width=0.95\textwidth]{slides-images/slide-104.jpg}
\caption{现代视频模型示意：token 数和 latent 尺度迅速上升，模型规模也随之暴涨\protect\footnotemark}
\end{figure}
\footnotetext{Slides 第105页。}

\begin{figure}[H]
\centering
\includegraphics[width=0.95\textwidth]{slides-images/slide-105.jpg}
\caption{视频 diffusion 的时代已经到来：多个大模型家族同时围绕视频生成展开竞争\protect\footnotemark}
\end{figure}
\footnotetext{Slides 第106页。}

\subsection{为什么还要 distillation}

Diffusion 的采样成本本来就高。latent diffusion 好不容易把每一步便宜化了，但如果还要跑 30--50 步，推理依旧会慢。于是 distillation 成为必需：把多步采样压成更少步，甚至一步，并且有时还能把 CFG 一起蒸馏进去。

\begin{figure}[H]
\centering
\includegraphics[width=0.95\textwidth]{slides-images/slide-107.jpg}
\caption{扩散蒸馏的动机：采样需要很多步，所以推理会很慢\protect\footnotemark}
\end{figure}
\footnotetext{Slides 第108页。}

\begin{importantbox}{蒸馏在这里解决的是什么}
蒸馏不改变“模型学什么”，主要是在改变“模型怎么被执行”：
\begin{itemize}
    \item 把多步采样压缩成少步
    \item 把 CFG 的效果直接 bake 进学生模型
    \item 让 diffusion 系统更接近产品级时延要求
\end{itemize}
\end{importantbox}

